\subsection{\texorpdfstring{取值于 \(\overline{R}\) 中的函数的等价类}{取值于 R-bar 中的函数的等价类}}

按照 No.~3 中的定义，我们说一个函数 \(f\) （在 \(X\) 中几乎处处有定义，取值于 \(\overline{\mathbf{R}}\) 中）几乎处处有限，如果由诸 \(x \in X\) （使 \(f(x)\) 有定义且有限）组成的集具有可忽略的余集。几乎处处有限的函数等价于一个处处有限的函数；因此可把它的类 \(\widetilde{f}\) 等同于定义在 \(X\) 上（或在 \(X\) 中几乎处处有定义）的有限数值函数的一个类。特别地，两个几乎处处有限的函数的类的和与积有定义，且这些类所成之集是 \(\mathbf{R}\) 上的一个代数。若 \((f_n)\) 是取值于 \(\overline{\mathbf{R}}\) 中、几乎处处有定义且有限的函数序列，则部分和 \(\sum_{k=1}^{n} f_k(x)\) 几乎处处有定义；若对几乎所有的 \(x \in X\) ，它们有极限 \(f(x)\) （在 \(\overline{R}\) 中），我们仍说以 \(f_{n}\) 为通项的级数几乎处处收敛，且 f 是该级数的和（注意 f 不一定几乎处处有限）。

若 \(f\) 与 \(g\) 是在 \(X\) 中几乎处处有定义且有限的两个数值函数，则 \(\widetilde{f} + \widetilde{g}\) （相应地 \(\widetilde{f}\widetilde{g}\) ）是每个等于 \(f(x) + g(x)\) （相应地 \(f(x)g(x)\) ）的函数的类（等式在使该表达式有意义的诸点 \(x \in X\) 处成立）。注意，\(f\) 与 \(g\) 可以都处处有定义，而 \(f(x) + g(x)\) （相应地 \(f(x)g(x)\) ）不对所有 \(x\) 有定义（GT, IV, §4, No.~3）；这时按定义 \(f + g\) （相应地 \(fg\) ）是在该表达式有定义的点处等于 \(f(x) + g(x)\) （相应地 \(f(x)g(x)\) ）的函数；因此它只是几乎处处有定义的。

设 \(f\) 与 \(g\) 是在 \(X\) 中几乎处处有定义的两个数值函数（有限或否），满足 \(f(x) \leqslant g(x)\) 几乎处处成立；若 \(f_1\) 等价于 \(f\) 且 \(g_1\) 等价于 \(g\) ，则显然 \(f_1(x) \leqslant g_1(x)\) 也几乎处处成立。因此所论关系只依赖于 \(f\) 与 \(g\) 的类；记作 \(\widetilde{f} \leqslant \widetilde{g}\) ，并且立即可以验证这个关系是取值于 \(\overline{\mathbf{R}}\) 中的函数的等价类所成之集上的一个序关系。若 \((\widetilde{f}_n)\) 是这样的一些类构成的可数族（有限或无穷），且对每个 \(n\) ，\(f_n\) 与 \(g_n\) 是两个几乎处处有定义、属于类 \(\widetilde{f}_n\) 的函数，则由 No.~4 的命题 8 推出，几乎处处有定义的函数 \(\sup_n f_n\) 与 \(\sup_n g_n\) 等价；因此它们的类只依赖于诸类 \(\widetilde{f}_n\) ，并且立即可验证它是这些类的上确界 \(\sup_n \widetilde{f}_n\) ——在按刚才所述方式排序的、取值于 \(\overline{\mathbf{R}}\) 中的函数的类所成之集中（从而这个集特别是一个格序集）。类似地可证明下确界 \(\inf_n \widetilde{f}_n\) 的存在性，并且有 \(\inf_n \widetilde{f}_n = -\sup_n (-\widetilde{f}_n)\) 。由此推出 \(\limsup_{n \to \infty} f_n\) 与 \(\limsup_{n \to \infty} g_n\) 也等价，它们的类记作 \(\limsup_{n \to \infty} \widetilde{f}_n\) ，等于 \(\inf_n (\sup_{p \geqslant 0} \widetilde{f}_{n+p})\) ；\(\liminf_{n \to \infty} \widetilde{f}_n\) 类似地定义。

数值函数 f（有限或否）如果等价于 0，就称为可忽略的；对于正的且处处有定义的函数，这个定义凭借定理 1 与定义 1 等价。要使 f 可忽略，必须且只须 \(|f|\) 可忽略（或 \(f^{+}\) 与 \(f^{-}\) 都可忽略）。

